
\section{Discussion}\label{sec:discussion}

\textbf{Classical vs.\ deep learning.}
SVM and RF on MFCC features outperform our replicated CNN (76.7\% vs.\ 60.4\%).
This does not imply that classical methods are generally better than deep learning for audio.
The CNN was trained without data augmentation (colored noise and impulse responses), which Fai\ss{} and Stowell~\cite{faiss2023adaptive} report as essential for achieving 82\%.
With augmentation, the CNN would likely outperform the classical models.

\textbf{Replication gap.}
Our replicated CNN achieves 60.4\% compared to 82\% in Fai\ss{} and Stowell~\cite{faiss2023adaptive}.
The primary cause is the absence of augmentation.
Secondary factors include training on a single random seed and slight differences in hardware (Apple MPS vs.\ CUDA).
We document this limitation rather than attempting to close the gap, as it falls outside the scope of this project.

\textbf{MFCC feature space.}
The NMI of 0.58 shows that MFCC statistics encode substantial species-discriminative information.
Low-order MFCC means, which capture the coarse spectral envelope, are the most informative features.
This aligns with the acoustic properties of insect calls, where species are distinguished primarily by call frequency, chirp rate, and syllable structure.

\textbf{Class imbalance.}
The dataset is imbalanced (10--160 recordings per species).
We use \texttt{class\_weight=balanced} for both SVM and RF to compensate.
Macro-F1 is a more reliable metric than accuracy in this setting because it weights each class equally.

\textbf{Interpretability trade-off.}
RF provides explicit feature importance while the CNN does not.
In ecological applications, understanding which acoustic cues drive a prediction is valuable for validating the model and building trust with domain experts.
The classical pipeline is therefore a practical complement to the deep learning baseline, not just a weaker alternative.
